\chapter{Design and Development of the Prolog Expert System}
\label{ch:design_and_development}

This chapter describes how the knowledge of Chapter~\ref{ch:theory_and_kr} was turned into a working expert system. It covers the architecture (Section~\ref{sec:architecture}), the design of the knowledge base (Section~\ref{sec:kb-design}), the predicates and rules (Section~\ref{sec:predicates}), the three inference mechanisms and their combination (Sections~\ref{sec:inference} to~\ref{sec:cbd}), the explanation facility (Section~\ref{sec:explanation}) and the corrective-maintenance planner (Section~\ref{sec:planner}). All diagnostic reasoning is implemented in SWI-Prolog. Python is used only to generate synthetic test data and to format results.

\lstset{language=Prolog, basicstyle=\ttfamily\footnotesize, columns=fullflexible, keepspaces=true,
        breaklines=true, frame=single, rulecolor=\color{black!35}, xleftmargin=2pt, xrightmargin=2pt,
        commentstyle=\color{black!60}\itshape, keywordstyle={}, showstringspaces=false}

% ---------------------------------------------------------------------------
\section{System Architecture}
\label{sec:architecture}

Figure~\ref{fig:architecture} shows the architecture. It follows the ten components that the project framework asks for:

\begin{enumerate}
  \item a user and input interface: test cases loaded as Prolog facts, and queries such as \texttt{why/1};
  \item observation and preprocessing predicates: noise estimation and graded symptoms;
  \item engineering-calculation predicates for the equations of Section~\ref{sec:maths};
  \item the knowledge base;
  \item a deductive inference engine;
  \item an abductive diagnosis module;
  \item a consistency-based diagnosis module;
  \item an explanation facility;
  \item a corrective-maintenance planner;
  \item an evaluation and reporting module.
\end{enumerate}

\begin{figure}[htbp]
  \centering
  % Figure 2: expert-system architecture.
\begin{tikzpicture}[x=1cm, y=1cm]
  \node[data=5cm] (obs) at (4.5,0) {Observations and test cases\\ \footnotesize\texttt{obs/4} facts in \texttt{data/}};
  \node[kb=5cm] (calc) at (4.5,-1.5) {Engineering calculations\\ \footnotesize Stage 4 equations, \texttt{kb/calculations.pl}};
  \node[reason=5cm] (evid) at (4.5,-3.0) {Noise estimation and graded symptoms\\ \footnotesize \texttt{engine/evidence.pl}};

  \node[kb=4.2cm, minimum height=3.4cm, align=left] (kb) at (12.1,-1.9)
    {\textbf{Knowledge base}\\[2pt]
     \footnotesize components and ratings\\
     \footnotesize thresholds with provenance\\
     \footnotesize 47 diagnostic rules\\
     \footnotesize causal signatures\\
     \footnotesize redundancy relations\\
     \footnotesize \quad and fault models\\
     \footnotesize maintenance plans};

  \node[reason=3.1cm] (ded) at (0.9,-4.9) {Deductive\\ \footnotesize rules, proof trees};
  \node[reason=3.1cm] (abd) at (4.5,-4.9) {Abductive\\ \footnotesize parsimonious covering};
  \node[reason=3.1cm] (cbd) at (8.1,-4.9) {Consistency-based\\ \footnotesize minimal hitting sets};

  \node[reason=5.4cm] (hyb) at (4.5,-6.6) {Hybrid combiner\\ \footnotesize abduction proposes, consistency disposes};
  \node[out=5cm] (diag) at (4.5,-8.1) {Diagnosis\\ \footnotesize set of fault labels};
  \node[out=3.6cm] (expl) at (2.3,-9.7) {Explanation facility\\ \footnotesize \texttt{why/1}};
  \node[out=3.6cm] (maint) at (6.7,-9.7) {Corrective-maintenance planner};
  \node[out=5cm] (eval) at (4.5,-11.3) {Evaluation and reporting\\ \footnotesize metrics, sensitivity, tables};

  \draw[flow] (obs) -- (calc);
  \draw[flow] (calc) -- (evid);
  \draw[flow] (calc.west) -| (ded.north);
  \draw[flow] (evid) -- (abd);
  \draw[flow] (evid.east) -| (cbd.north);
  \draw[flow] (ded.south) |- (hyb.west);
  \draw[flow] (abd) -- (hyb);
  \draw[flow] (cbd.south) |- (hyb.east);
  \draw[flow] (hyb) -- (diag);
  \draw[flow] (diag.south) -- ++(0,-0.35) -| (expl.north);
  \draw[flow] (diag.south) -- ++(0,-0.35) -| (maint.north);
  \draw[flow] (expl.south) -- (expl.south |- eval.north);
  \draw[flow] (maint.south) -- (maint.south |- eval.north);

  \draw[uses] (kb.west |- calc.east) -- (calc.east);
  \draw[uses] (kb.west |- evid.east) -- (evid.east);
  \draw[uses] (kb.south) |- (cbd.east);
  \node[tag, anchor=south] at (9.0,-1.5) {uses};
\end{tikzpicture}

  \caption{Architecture of the Prolog expert system.}
  \label{fig:architecture}
\end{figure}

The knowledge base is kept separate from the engines that use it, which is the defining property of an expert system \parencite{jackson1998}. Table~\ref{tab:modules} lists the source files. The knowledge base and engines together are about 1{,}900 lines of Prolog, and the test and reporting code about 460 more. A single entry file, \texttt{main.pl}, loads everything.

\begin{table}[htbp]
  \centering
  \small
  \caption{Source modules.}
  \label{tab:modules}
  \begin{tabular}{L{4.8cm}L{9.2cm}}
    \toprule
    File & Role \\
    \midrule
    \texttt{kb/components.pl} & Components, topology, ratings and the variable dictionary \\
    \texttt{kb/thresholds.pl}, \texttt{kb/sources.pl} & Reference values and their provenance \\
    \texttt{kb/calculations.pl} & Engineering equations as predicates \\
    \texttt{kb/rules.pl} & Diagnostic rules stored as \texttt{drule/4} facts \\
    \texttt{kb/symptoms.pl} & Graded symptom definitions with noise models \\
    \texttt{kb/causal.pl} & Causal signatures for abduction \\
    \texttt{kb/arr.pl} & Redundancy relations and fault models \\
    \texttt{kb/maintenance.pl} & Corrective-maintenance plans and plan interactions \\
    \texttt{engine/prover.pl} & Proof-recording meta-interpreter \\
    \texttt{engine/deductive.pl} & Deductive diagnosis and closed-world conclusions \\
    \texttt{engine/evidence.pl} & Noise estimation and evidence grading \\
    \texttt{engine/abductive.pl} & Parsimonious covering \\
    \texttt{engine/consistency.pl} & Conflicts, minimal hitting sets and fault models \\
    \texttt{engine/hybrid.pl} & The four reasoning configurations \\
    \texttt{engine/explain.pl} & The explanation facility \\
    \texttt{tests/*.pl} & Test runner, metrics, sensitivity and table export \\
    \bottomrule
  \end{tabular}
\end{table}

% ---------------------------------------------------------------------------
\section{Knowledge-Base Design}
\label{sec:kb-design}

The knowledge base follows the seven layers of Table~\ref{tab:kr-layers}. Its contents are summarised in Table~\ref{tab:kb-size}.

\begin{table}[htbp]
  \centering
  \small
  \caption{Contents of the knowledge base.}
  \label{tab:kb-size}
  \begin{tabular}{L{7.0cm}L{2.0cm}L{5.0cm}}
    \toprule
    Element & Count & Representation \\
    \midrule
    Diagnostic rules deriving conditions & 29 & \texttt{drule/4} \\
    Diagnostic rules deriving faults & 18 & \texttt{drule/4} \\
    Closed-world conclusions & 2 & \texttt{engine/deductive.pl} \\
    Thresholds & 24 & \texttt{threshold/4} \\
    Sources: standard, literature, method, assumption & 3, 1, 3, 17 & \texttt{source/3} \\
    Engineering-calculation predicates & 20 & \texttt{calc\_pred/1} \\
    Graded, yes/no and composite symptoms & 20, 2, 2 & \texttt{sdef/7} and others \\
    Causal evidence links & 31 & \texttt{evidence/3} \\
    Required symptoms, masking links, extra explanations & 12, 2, 4 & \texttt{requires/2} and others \\
    Analytical redundancy relations & 20 & \texttt{arr/3} \\
    Fault models & 14 & \texttt{fm/2} \\
    Maintenance plans & 12 & \texttt{maintenance/5} \\
    Plan interactions & 4 & \texttt{plan\_interaction/3} \\
    \bottomrule
  \end{tabular}
\end{table}

\subsection{Facts}

Structural knowledge is stated as ground facts, for example \texttt{component(crah\_1, crah)}, \texttt{serves(crah\_1, rack\_a04)} and \texttt{rated(crah\_1, fan\_power\_rated, 5.5, kw)}. Observations are facts \texttt{obs(Case, Variable, Entity, Value)}, so every case is a self-contained set of facts and the same knowledge base diagnoses all cases.

\subsection{Thresholds and provenance}

Each threshold is a fact with its value, unit and source, and each source has a type:
\begin{lstlisting}
threshold(t_inlet_rec_max, 27.0, degC, ashrae_tc99_2021).
threshold(strat_out_max,    3.5, kelvin, a7_outlet_strat).
source(a7_outlet_strat, assumption,
       'Engineering assumption: a top-to-bottom outlet difference above 3.5 K
        confirms recirculation. Normal racks stay below 3 K.').
\end{lstlisting}
Rules and symptoms refer to thresholds by name through \texttt{threshold/2}, never by value. A threshold can therefore be changed in one place, and the sensitivity analysis of Chapter~\ref{ch:dataset_and_evaluation} can perturb it without touching any rule. The complete threshold and provenance tables are generated from these facts (Tables~\ref{tab:thresholds} and~\ref{tab:provenance}).

\subsection{Engineering calculations}

Each equation of Section~\ref{sec:maths} that the reasoning needs is a Horn clause with arithmetic, commented with its equation tag. For example, the fan affinity law and the fan power ratio, Equations E23 and E24, are:
\begin{lstlisting}
%   Eq. 23  P_f,exp = P_f,rated * (N / N_rated)^3      fan affinity law
expected_fan_power(C, U, P) :-
    obs(C, fan_speed, U, N),
    rated(U, fan_power_rated, Pr, _),
    rated(U, fan_speed_rated, Nr, _),
    P is Pr * (N / Nr) ** 3.

%   Eq. 24  rho_f = P_f,measured / P_f,exp
fan_power_ratio(C, U, Ratio) :-
    obs(C, fan_power, U, Pm),
    expected_fan_power(C, U, Pe),
    Pe > 0,
    Ratio is Pm / Pe.
\end{lstlisting}
A calculation fails when one of its inputs is missing, and a rule that uses it then cannot fire. This is how incomplete evidence is handled in the deductive mode.

% ---------------------------------------------------------------------------
\section{Predicates and Rules}
\label{sec:predicates}

Table~\ref{tab:predicates} lists the main predicates. The diagnostic rules are listed in full in Table~\ref{tab:rules}, which is generated from \texttt{kb/rules.pl}. Each rule is a Horn clause stored as data:
\begin{lstlisting}
%   R24  F4 from exhaust evidence alone, when the top inlet sensor is
%        suspect because it reads low.
drule(r24, fault(C, hot_air_recirculation(R)),
      [ cond(C, outlet_stratified(R)), cond(C, inlet_not_stratified(R)),
        serves(U, R), cond(C, supply_air_normal(U)) ],
      eu_coc_dc).
\end{lstlisting}
It reads: rack $R$ has hot-air recirculation if its exhaust is stratified, its inlet reading is not, and its cooling unit supplies normal air. The rule exists because a low-reading top sensor can hide real recirculation (Table~\ref{tab:interactions}). The exhaust then carries the evidence that the inlet sensor does not.

\begin{table}[htbp]
  \centering
  \small
  \caption{Main predicates.}
  \label{tab:predicates}
  \begin{tabular}{L{4.6cm}L{9.4cm}}
    \toprule
    Predicate & Purpose \\
    \midrule
    \texttt{obs(C, V, E, X)} & Observation: variable $V$ of entity $E$ in case $C$ has value $X$ \\
    \texttt{threshold(N, X)} & Value of threshold $N$, after any sensitivity perturbation \\
    \texttt{drule(Id, H, B, S)} & Diagnostic rule $H \leftarrow B$ with identifier and source \\
    \texttt{prove(G, P)} & Proves goal $G$ and returns its proof tree $P$ \\
    \texttt{diagnose(C, Ds)} & Deductive diagnosis with proofs \\
    \texttt{case\_symptoms(C, Ss)} & Graded symptoms with evidence and status \\
    \texttt{abduce(C, R)} & Abductive candidates, minimal explanations and verdicts \\
    \texttt{conflicts(Ss, Ks)} & Conflict sets from violated redundancy relations \\
    \texttt{cbd\_diagnoses(Ss, Ds)} & Minimal consistent diagnoses \\
    \texttt{diagnose\_mode(M, C, Ls)} & Diagnosis in configuration $M$ \\
    \texttt{why(C)} & Full explanation of the complete system's diagnosis \\
    \texttt{maintenance(K, Pr, As, R, S)} & Plan for fault class $K$: priority, actions, rationale, source \\
    \bottomrule
  \end{tabular}
\end{table}

% Generated by tests/export_latex.pl from the knowledge base. Do not edit by hand.
\begin{footnotesize}
\begin{longtable}{L{0.9cm}L{4.0cm}L{7.0cm}L{1.9cm}}
  \caption{Diagnostic rules, generated from \texttt{kb/rules.pl}. In each rule \texttt{A} is the case being diagnosed and \texttt{\textbackslash+} is negation as failure.}\label{tab:rules}\\
  \toprule
  Rule & Concludes & If all of these hold & Source type \\
  \midrule
  \endfirsthead
  \toprule
  Rule & Concludes & If all of these hold & Source type \\
  \midrule
  \endhead
  \texttt{r01} & condition \texttt{rack\_\allowbreak{}inlet\_\allowbreak{}hot(B)} & \texttt{component(B,\allowbreak{} rack), \allowbreak inlet\_\allowbreak{}max(A,\allowbreak{} B,\allowbreak{} C), \allowbreak threshold(t\_\allowbreak{}inlet\_\allowbreak{}rec\_\allowbreak{}max,\allowbreak{} D), \allowbreak C>D} & standard \\
  \texttt{r02} & condition \texttt{row\_\allowbreak{}inlet\_\allowbreak{}hot} & \texttt{hot\_\allowbreak{}rack\_\allowbreak{}fraction(A,\allowbreak{} B), \allowbreak threshold(row\_\allowbreak{}hot\_\allowbreak{}fraction,\allowbreak{} C), \allowbreak B>=C} & assumption \\
  \texttt{r03} & condition \texttt{localised\_\allowbreak{}inlet\_\allowbreak{}hot(B)} & \texttt{cond(A,\allowbreak{} rack\_\allowbreak{}inlet\_\allowbreak{}hot(B)), \allowbreak hot\_\allowbreak{}rack\_\allowbreak{}fraction(A,\allowbreak{} C), \allowbreak threshold(row\_\allowbreak{}hot\_\allowbreak{}fraction,\allowbreak{} D), \allowbreak C<D} & assumption \\
  \texttt{r04} & condition \texttt{inlet\_\allowbreak{}stratified(B)} & \texttt{inlet\_\allowbreak{}stratification(A,\allowbreak{} B,\allowbreak{} C), \allowbreak threshold(strat\_\allowbreak{}max,\allowbreak{} D), \allowbreak C>D} & assumption \\
  \texttt{r05} & condition \texttt{supply\_\allowbreak{}air\_\allowbreak{}hot(B)} & \texttt{component(B,\allowbreak{} crah), \allowbreak obs(A,\allowbreak{} t\_\allowbreak{}supply\_\allowbreak{}air,\allowbreak{} B,\allowbreak{} C), \allowbreak threshold(t\_\allowbreak{}supply\_\allowbreak{}setpoint,\allowbreak{} D), \allowbreak threshold(t\_\allowbreak{}supply\_\allowbreak{}band,\allowbreak{} E), \allowbreak C>D+E} & assumption \\
  \texttt{r06} & condition \texttt{supply\_\allowbreak{}air\_\allowbreak{}normal(B)} & \texttt{component(B,\allowbreak{} crah), \allowbreak obs(A,\allowbreak{} t\_\allowbreak{}supply\_\allowbreak{}air,\allowbreak{} B,\allowbreak{} C), \allowbreak threshold(t\_\allowbreak{}supply\_\allowbreak{}setpoint,\allowbreak{} D), \allowbreak threshold(t\_\allowbreak{}supply\_\allowbreak{}band,\allowbreak{} E), \allowbreak C=<D+E} & assumption \\
  \texttt{r07} & condition \texttt{fan\_\allowbreak{}power\_\allowbreak{}low(B)} & \texttt{component(B,\allowbreak{} crah), \allowbreak fan\_\allowbreak{}power\_\allowbreak{}ratio(A,\allowbreak{} B,\allowbreak{} C), \allowbreak threshold(fan\_\allowbreak{}ratio\_\allowbreak{}min,\allowbreak{} D), \allowbreak C<D} & assumption \\
  \texttt{r08} & condition \texttt{fan\_\allowbreak{}power\_\allowbreak{}normal(B)} & \texttt{component(B,\allowbreak{} crah), \allowbreak fan\_\allowbreak{}power\_\allowbreak{}ratio(A,\allowbreak{} B,\allowbreak{} C), \allowbreak threshold(fan\_\allowbreak{}ratio\_\allowbreak{}min,\allowbreak{} D), \allowbreak C>=D} & literature \\
  \texttt{r09} & condition \texttt{valve\_\allowbreak{}open(B)} & \texttt{obs(A,\allowbreak{} valve\_\allowbreak{}position,\allowbreak{} B,\allowbreak{} C), \allowbreak threshold(valve\_\allowbreak{}open\_\allowbreak{}min,\allowbreak{} D), \allowbreak C>=D} & assumption \\
  \texttt{r10} & condition \texttt{chw\_\allowbreak{}flow\_\allowbreak{}low(B)} & \texttt{chw\_\allowbreak{}flow\_\allowbreak{}ratio(A,\allowbreak{} B,\allowbreak{} C), \allowbreak threshold(chw\_\allowbreak{}flow\_\allowbreak{}ratio\_\allowbreak{}min,\allowbreak{} D), \allowbreak C<D} & assumption \\
  \texttt{r20} & fault \texttt{crah\_\allowbreak{}fan\_\allowbreak{}failure(B)} & \texttt{cond(A,\allowbreak{} row\_\allowbreak{}inlet\_\allowbreak{}hot), \allowbreak cond(A,\allowbreak{} fan\_\allowbreak{}power\_\allowbreak{}low(B))} & literature \\
  \texttt{r21} & fault \texttt{chw\_\allowbreak{}supply\_\allowbreak{}loss(B)} & \texttt{cond(A,\allowbreak{} supply\_\allowbreak{}air\_\allowbreak{}hot(B)), \allowbreak cond(A,\allowbreak{} valve\_\allowbreak{}open(B)), \allowbreak cond(A,\allowbreak{} chw\_\allowbreak{}flow\_\allowbreak{}low(B)), \allowbreak cond(A,\allowbreak{} fan\_\allowbreak{}power\_\allowbreak{}normal(B))} & assumption \\
  \texttt{r11} & condition \texttt{outlet\_\allowbreak{}stratified(B)} & \texttt{outlet\_\allowbreak{}stratification(A,\allowbreak{} B,\allowbreak{} C), \allowbreak threshold(strat\_\allowbreak{}out\_\allowbreak{}max,\allowbreak{} D), \allowbreak C>D} & assumption \\
  \texttt{r12} & condition \texttt{outlet\_\allowbreak{}not\_\allowbreak{}stratified(B)} & \texttt{outlet\_\allowbreak{}stratification(A,\allowbreak{} B,\allowbreak{} C), \allowbreak threshold(strat\_\allowbreak{}out\_\allowbreak{}max,\allowbreak{} D), \allowbreak C=<D} & assumption \\
  \texttt{r13} & condition \texttt{inlet\_\allowbreak{}not\_\allowbreak{}stratified(B)} & \texttt{inlet\_\allowbreak{}stratification(A,\allowbreak{} B,\allowbreak{} C), \allowbreak threshold(strat\_\allowbreak{}max,\allowbreak{} D), \allowbreak C=<D} & assumption \\
  \texttt{r14} & condition \texttt{filter\_\allowbreak{}loaded(B)} & \texttt{component(B,\allowbreak{} crah), \allowbreak filter\_\allowbreak{}loading(A,\allowbreak{} B,\allowbreak{} C), \allowbreak threshold(filter\_\allowbreak{}loading\_\allowbreak{}max,\allowbreak{} D), \allowbreak C>D} & assumption \\
  \texttt{r15} & condition \texttt{server\_\allowbreak{}fan\_\allowbreak{}stopped(B)} & \texttt{obs(A,\allowbreak{} server\_\allowbreak{}fan\_\allowbreak{}rpm,\allowbreak{} B,\allowbreak{} C), \allowbreak threshold(server\_\allowbreak{}fan\_\allowbreak{}min,\allowbreak{} D), \allowbreak C<D} & assumption \\
  \texttt{r16} & condition \texttt{cpu\_\allowbreak{}hot(B)} & \texttt{obs(A,\allowbreak{} t\_\allowbreak{}cpu,\allowbreak{} B,\allowbreak{} C), \allowbreak threshold(cpu\_\allowbreak{}hot\_\allowbreak{}min,\allowbreak{} D), \allowbreak C>=D} & assumption \\
  \texttt{r17} & condition \texttt{branch\_\allowbreak{}overloaded(B)} & \texttt{branch\_\allowbreak{}loading(A,\allowbreak{} B,\allowbreak{} C), \allowbreak threshold(branch\_\allowbreak{}loading\_\allowbreak{}max,\allowbreak{} D), \allowbreak C>D} & standard \\
  \texttt{r18} & condition \texttt{cell\_\allowbreak{}resistance\_\allowbreak{}high(B)} & \texttt{component(B,\allowbreak{} battery), \allowbreak cell\_\allowbreak{}resistance\_\allowbreak{}ratio(A,\allowbreak{} B,\allowbreak{} C), \allowbreak threshold(cell\_\allowbreak{}r\_\allowbreak{}ratio\_\allowbreak{}max,\allowbreak{} D), \allowbreak C>D} & assumption \\
  \texttt{r19} & condition \texttt{ups\_\allowbreak{}efficiency\_\allowbreak{}bypass\_\allowbreak{}like(B)} & \texttt{component(B,\allowbreak{} ups), \allowbreak ups\_\allowbreak{}efficiency(A,\allowbreak{} B,\allowbreak{} C), \allowbreak threshold(ups\_\allowbreak{}eff\_\allowbreak{}bypass\_\allowbreak{}min,\allowbreak{} D), \allowbreak C>=D} & assumption \\
  \texttt{r40} & condition \texttt{mid\_\allowbreak{}inlet\_\allowbreak{}deviates(B)} & \texttt{component(B,\allowbreak{} rack), \allowbreak mid\_\allowbreak{}inlet\_\allowbreak{}deviation(A,\allowbreak{} B,\allowbreak{} C), \allowbreak threshold(neighbour\_\allowbreak{}dev\_\allowbreak{}max,\allowbreak{} D), \allowbreak abs(C)>D, \allowbreak \textbackslash{}+cond(A,\allowbreak{} inlet\_\allowbreak{}rise\_\allowbreak{}corroborated(B))} & assumption \\
  \texttt{r45} & condition \texttt{inlet\_\allowbreak{}rise\_\allowbreak{}corroborated(B)} & \texttt{mid\_\allowbreak{}inlet\_\allowbreak{}deviation(A,\allowbreak{} B,\allowbreak{} C), \allowbreak top\_\allowbreak{}inlet\_\allowbreak{}deviation(A,\allowbreak{} B,\allowbreak{} D), \allowbreak C*D>0, \allowbreak threshold(neighbour\_\allowbreak{}dev\_\allowbreak{}max,\allowbreak{} E), \allowbreak abs(D)>E} & assumption \\
  \texttt{r46} & condition \texttt{top\_\allowbreak{}inlet\_\allowbreak{}sensor\_\allowbreak{}suspect(B)} & \texttt{component(B,\allowbreak{} rack), \allowbreak inlet\_\allowbreak{}stratification(A,\allowbreak{} B,\allowbreak{} C), \allowbreak threshold(neighbour\_\allowbreak{}dev\_\allowbreak{}max,\allowbreak{} D), \allowbreak C< -D, \allowbreak \textbackslash{}+cond(A,\allowbreak{} mid\_\allowbreak{}inlet\_\allowbreak{}deviates(B))} & method \\
  \texttt{r41} & condition \texttt{top\_\allowbreak{}inlet\_\allowbreak{}sensor\_\allowbreak{}suspect(B)} & \texttt{cond(A,\allowbreak{} inlet\_\allowbreak{}stratified(B)), \allowbreak cond(A,\allowbreak{} outlet\_\allowbreak{}not\_\allowbreak{}stratified(B)), \allowbreak \textbackslash{}+cond(A,\allowbreak{} mid\_\allowbreak{}inlet\_\allowbreak{}deviates(B))} & assumption \\
  \texttt{r42} & condition \texttt{top\_\allowbreak{}inlet\_\allowbreak{}sensor\_\allowbreak{}suspect(B)} & \texttt{cond(A,\allowbreak{} outlet\_\allowbreak{}stratified(B)), \allowbreak cond(A,\allowbreak{} inlet\_\allowbreak{}not\_\allowbreak{}stratified(B)), \allowbreak \textbackslash{}+cond(A,\allowbreak{} mid\_\allowbreak{}inlet\_\allowbreak{}deviates(B))} & assumption \\
  \texttt{r43} & condition \texttt{supply\_\allowbreak{}sensor\_\allowbreak{}inconsistent(B)} & \texttt{component(B,\allowbreak{} crah), \allowbreak cond(A,\allowbreak{} fan\_\allowbreak{}power\_\allowbreak{}nominal(B)), \allowbreak supply\_\allowbreak{}air\_\allowbreak{}deviation(A,\allowbreak{} B,\allowbreak{} C), \allowbreak threshold(neighbour\_\allowbreak{}dev\_\allowbreak{}max,\allowbreak{} D), \allowbreak abs(C)>D, \allowbreak coil\_\allowbreak{}balance\_\allowbreak{}residual(A,\allowbreak{} B,\allowbreak{} E), \allowbreak threshold(balance\_\allowbreak{}resid\_\allowbreak{}max,\allowbreak{} F), \allowbreak abs(E)>F, \allowbreak coil\_\allowbreak{}balance\_\allowbreak{}residual\_\allowbreak{}corrected(A,\allowbreak{} B,\allowbreak{} G), \allowbreak abs(G)=<F} & method \\
  \texttt{r48} & condition \texttt{fan\_\allowbreak{}power\_\allowbreak{}nominal(B)} & \texttt{component(B,\allowbreak{} crah), \allowbreak fan\_\allowbreak{}power\_\allowbreak{}ratio(A,\allowbreak{} B,\allowbreak{} C), \allowbreak threshold(fan\_\allowbreak{}ratio\_\allowbreak{}reduced,\allowbreak{} D), \allowbreak C>=D} & assumption \\
  \texttt{r49} & condition \texttt{fan\_\allowbreak{}power\_\allowbreak{}reduced(B)} & \texttt{component(B,\allowbreak{} crah), \allowbreak fan\_\allowbreak{}power\_\allowbreak{}ratio(A,\allowbreak{} B,\allowbreak{} C), \allowbreak threshold(fan\_\allowbreak{}ratio\_\allowbreak{}reduced,\allowbreak{} D), \allowbreak C<D} & assumption \\
  \texttt{r47} & condition \texttt{supply\_\allowbreak{}air\_\allowbreak{}cold(B)} & \texttt{component(B,\allowbreak{} crah), \allowbreak obs(A,\allowbreak{} t\_\allowbreak{}supply\_\allowbreak{}air,\allowbreak{} B,\allowbreak{} C), \allowbreak threshold(t\_\allowbreak{}supply\_\allowbreak{}setpoint,\allowbreak{} D), \allowbreak threshold(t\_\allowbreak{}supply\_\allowbreak{}band,\allowbreak{} E), \allowbreak C<D-E} & assumption \\
  \texttt{r44} & condition \texttt{telemetry\_\allowbreak{}stale(B)} & \texttt{obs(A,\allowbreak{} sample\_\allowbreak{}age,\allowbreak{} B,\allowbreak{} C), \allowbreak threshold(sample\_\allowbreak{}age\_\allowbreak{}max,\allowbreak{} D), \allowbreak C>D} & assumption \\
  \texttt{r36} & fault \texttt{crah\_\allowbreak{}fan\_\allowbreak{}failure(B)} & \texttt{cond(A,\allowbreak{} row\_\allowbreak{}inlet\_\allowbreak{}hot), \allowbreak cond(A,\allowbreak{} supply\_\allowbreak{}air\_\allowbreak{}cold(B)), \allowbreak \textbackslash{}+cond(A,\allowbreak{} supply\_\allowbreak{}sensor\_\allowbreak{}inconsistent(B))} & assumption \\
  \texttt{r38} & fault \texttt{crah\_\allowbreak{}fan\_\allowbreak{}failure(B)} & \texttt{cond(A,\allowbreak{} fan\_\allowbreak{}power\_\allowbreak{}reduced(B)), \allowbreak cond(A,\allowbreak{} row\_\allowbreak{}inlet\_\allowbreak{}hot)} & assumption \\
  \texttt{r37} & fault \texttt{chw\_\allowbreak{}supply\_\allowbreak{}loss(B)} & \texttt{cond(A,\allowbreak{} supply\_\allowbreak{}air\_\allowbreak{}hot(B)), \allowbreak cond(A,\allowbreak{} valve\_\allowbreak{}open(B)), \allowbreak cond(A,\allowbreak{} fan\_\allowbreak{}power\_\allowbreak{}normal(B)), \allowbreak \textbackslash{}+obs(A,\allowbreak{} m\_\allowbreak{}chw,\allowbreak{} B,\allowbreak{} C), \allowbreak \textbackslash{}+cond(A,\allowbreak{} supply\_\allowbreak{}sensor\_\allowbreak{}inconsistent(B))} & assumption \\
  \texttt{r22} & fault \texttt{hot\_\allowbreak{}air\_\allowbreak{}recirculation(B)} & \texttt{cond(A,\allowbreak{} localised\_\allowbreak{}inlet\_\allowbreak{}hot(B)), \allowbreak serves(C,\allowbreak{} B), \allowbreak cond(A,\allowbreak{} supply\_\allowbreak{}air\_\allowbreak{}normal(C)), \allowbreak cond(A,\allowbreak{} inlet\_\allowbreak{}stratified(B)), \allowbreak \textbackslash{}+cond(A,\allowbreak{} top\_\allowbreak{}inlet\_\allowbreak{}sensor\_\allowbreak{}suspect(B))} & standard \\
  \texttt{r23} & fault \texttt{hot\_\allowbreak{}air\_\allowbreak{}recirculation(B)} & \texttt{cond(A,\allowbreak{} inlet\_\allowbreak{}stratified(B)), \allowbreak cond(A,\allowbreak{} outlet\_\allowbreak{}stratified(B)), \allowbreak serves(C,\allowbreak{} B), \allowbreak cond(A,\allowbreak{} supply\_\allowbreak{}air\_\allowbreak{}normal(C))} & standard \\
  \texttt{r24} & fault \texttt{hot\_\allowbreak{}air\_\allowbreak{}recirculation(B)} & \texttt{cond(A,\allowbreak{} outlet\_\allowbreak{}stratified(B)), \allowbreak cond(A,\allowbreak{} inlet\_\allowbreak{}not\_\allowbreak{}stratified(B)), \allowbreak serves(C,\allowbreak{} B), \allowbreak cond(A,\allowbreak{} supply\_\allowbreak{}air\_\allowbreak{}normal(C))} & standard \\
  \texttt{r25} & fault \texttt{filter\_\allowbreak{}clogging(B)} & \texttt{cond(A,\allowbreak{} filter\_\allowbreak{}loaded(B)), \allowbreak cond(A,\allowbreak{} fan\_\allowbreak{}power\_\allowbreak{}normal(B))} & assumption \\
  \texttt{r26} & fault \texttt{server\_\allowbreak{}fan\_\allowbreak{}failure(B)} & \texttt{cond(A,\allowbreak{} server\_\allowbreak{}fan\_\allowbreak{}stopped(B)), \allowbreak cond(A,\allowbreak{} cpu\_\allowbreak{}hot(B))} & assumption \\
  \texttt{r27} & fault \texttt{branch\_\allowbreak{}overload(B)} & \texttt{cond(A,\allowbreak{} branch\_\allowbreak{}overloaded(B))} & standard \\
  \texttt{r28} & fault \texttt{battery\_\allowbreak{}degradation(B)} & \texttt{cond(A,\allowbreak{} cell\_\allowbreak{}resistance\_\allowbreak{}high(B))} & assumption \\
  \texttt{r29} & fault \texttt{ups\_\allowbreak{}on\_\allowbreak{}bypass(B)} & \texttt{component(B,\allowbreak{} ups), \allowbreak obs(A,\allowbreak{} ups\_\allowbreak{}mode,\allowbreak{} B,\allowbreak{} bypass)} & assumption \\
  \texttt{r35} & fault \texttt{ups\_\allowbreak{}on\_\allowbreak{}bypass(B)} & \texttt{component(B,\allowbreak{} ups), \allowbreak \textbackslash{}+obs(A,\allowbreak{} ups\_\allowbreak{}mode,\allowbreak{} B,\allowbreak{} C), \allowbreak cond(A,\allowbreak{} ups\_\allowbreak{}efficiency\_\allowbreak{}bypass\_\allowbreak{}like(B))} & assumption \\
  \texttt{r50} & fault \texttt{sensor\_\allowbreak{}bias(t\_\allowbreak{}inlet\_\allowbreak{}mid,\allowbreak{} B)} & \texttt{cond(A,\allowbreak{} mid\_\allowbreak{}inlet\_\allowbreak{}deviates(B))} & assumption \\
  \texttt{r51} & fault \texttt{sensor\_\allowbreak{}bias(t\_\allowbreak{}inlet\_\allowbreak{}top,\allowbreak{} B)} & \texttt{cond(A,\allowbreak{} top\_\allowbreak{}inlet\_\allowbreak{}sensor\_\allowbreak{}suspect(B))} & assumption \\
  \texttt{r52} & fault \texttt{sensor\_\allowbreak{}bias(t\_\allowbreak{}supply\_\allowbreak{}air,\allowbreak{} B)} & \texttt{cond(A,\allowbreak{} supply\_\allowbreak{}sensor\_\allowbreak{}inconsistent(B))} & method \\
  \texttt{r53} & fault \texttt{telemetry\_\allowbreak{}dropout(B)} & \texttt{cond(A,\allowbreak{} telemetry\_\allowbreak{}stale(B))} & assumption \\
  \texttt{r30} & fault \texttt{normal} & no fault rule fires, and every rack inlet is within the recommended envelope & method \\
  \texttt{r31} & \texttt{unexplained\_\allowbreak{}abnormality} & no fault rule fires, but some rack inlet is outside the envelope & method \\
  \bottomrule
\end{longtable}
\end{footnotesize}


% ---------------------------------------------------------------------------
\section{Inference Mechanisms}
\label{sec:inference}

Figure~\ref{fig:diagnostic} shows how the three inference mechanisms work together on one case. The readings are first used to estimate measurement noise. Every symptom is then graded. Deductive rules, abduction and consistency checking run on the result, and the hybrid combiner chooses the final diagnosis.

\begin{figure}[htbp]
  \centering
  % Figure 5: diagnostic reasoning workflow of the complete system.
\begin{tikzpicture}[x=1cm, y=1cm]
  \node[data=6cm] (obs) at (4.5,0) {Readings of one case};
  \node[reason=6cm] (noise) at (4.5,-1.3) {Estimate noise from redundant sensors\\ \footnotesize $\sigma_T$ from middle inlets, $\kappa_P$ from rack power balances};
  \node[reason=6cm] (grade) at (4.5,-2.75) {Grade every symptom\\ \footnotesize present, absent, uncertain or unknown};

  \node[reason=3.6cm] (ded) at (0.3,-4.6) {Deductive rules\\ \footnotesize prove faults, record proof trees};
  \node[reason=3.6cm] (abd) at (4.5,-4.6) {Abduction\\ \footnotesize minimal consistent covers of present symptoms};
  \node[reason=3.6cm] (cbd) at (8.7,-4.6) {Consistency\\ \footnotesize violated relations give conflict sets};

  \node[reason=3.6cm] (rank) at (4.5,-6.4) {Rank covers by confidence\\ \footnotesize bonus if rules also prove a fault};
  \node[decision=2.4cm] (dec) at (4.5,-9.0) {next cover hits every conflict and respects fault models?};
  \node[reason=3.6cm] (mhs) at (10.3,-9.0) {Minimal hitting sets with fault models};
  \node[out=3.6cm] (yes) at (4.5,-11.3) {Diagnosis: that cover};
  \node[out=3.6cm] (no) at (10.3,-11.3) {Diagnosis: closest minimal consistent set};

  \draw[flow] (obs) -- (noise);
  \draw[flow] (noise) -- (grade);
  \draw[flow] (grade.south) -- ++(0,-0.4) -| (ded.north);
  \draw[flow] (grade) -- (abd);
  \draw[flow] (grade.south) -- ++(0,-0.4) -| (cbd.north);
  \draw[flow] (abd) -- (rank);
  \draw[uses] (ded.south) |- node[tag, pos=0.75, above] {bonus} (rank.west);
  \draw[flow] (rank) -- (dec);
  \draw[flow] (dec) -- node[tag, right] {yes} (yes);
  \draw[flow] (dec) -- node[tag, above, pos=0.45] {none left} (mhs);
  \draw[flow] (cbd.south) -- node[tag, right] {conflict sets} (cbd.south |- mhs.north);
  \draw[flow] (mhs) -- (no);
\end{tikzpicture}

  \caption{Diagnostic reasoning workflow of the complete system: abduction proposes, consistency disposes.}
  \label{fig:diagnostic}
\end{figure}

\subsection{Graded, noise-adaptive evidence}

For each symptom definition, \texttt{engine/evidence.pl} computes the feature value, compares it with its threshold and grades the evidence with Equation E55. The noise $\sigma_j$ of each feature is propagated from the datasheet accuracy of its sensors and scaled by the noise level estimated from the case. For temperature features this is the robust spread of the six middle inlets, and for electrical features the spread of the six rack power balances. A symptom is \emph{present} once its value is one standard deviation past the threshold, \emph{absent} once it is one standard deviation on the normal side, and \emph{uncertain} in between. If an input is missing the symptom is \emph{unknown}, which is never treated as absent. The graded symptoms of a case are computed once and cached, so that the abductive and consistency-based modules see exactly the same evidence.

\subsection{The four configurations}

The framework asks for four configurations to be compared, and \texttt{engine/hybrid.pl} provides them through \texttt{diagnose\_mode/3}:

\begin{itemize}
  \item \textbf{Deductive only:} the rules of Table~\ref{tab:rules}.
  \item \textbf{Deductive and abductive:} the best parsimonious explanation, with a bonus of 0.1 to the confidence of hypotheses that the rules also prove.
  \item \textbf{Deductive and consistency-based:} the deductive faults are kept if they form a consistent diagnosis, after dropping any fault whose component hits no conflict. Otherwise the minimal consistent diagnosis closest to them replaces them.
  \item \textbf{Complete system:} abduction proposes and consistency disposes. The best-ranked abductive explanation that is also a consistent diagnosis is accepted. If none is consistent, the minimal consistent diagnosis closest to the top explanation is used.
\end{itemize}
When several consistent diagnoses tie, the one with fewer sensor components wins, because assuming a physical fault is the safer maintenance choice.

% ---------------------------------------------------------------------------
\section{Deductive Reasoning}
\label{sec:deductive}

Deductive diagnosis implements Equation~\eqref{eq:deductive}. Because rules are stored as data, a small meta-interpreter is needed to prove them \parencite{bratko2012}. It returns a proof tree for every conclusion:
\begin{lstlisting}
derived(cond(_, _)).
derived(fault(_, _)).

prove(\+ G, naf(G)) :-                 % negation as failure
    !,
    \+ prove(G, _).
prove(G, by(Id, G, SubProofs)) :-      % a derived goal: try each rule
    derived(G),
    !,
    drule(Id, G, Body, _),
    prove_body(Body, SubProofs).
prove(G, leaf(G)) :-                   % a fact, calculation or comparison
    call(G).

prove_body([], []).
prove_body([G|Gs], [P|Ps]) :-
    prove(G, P),
    prove_body(Gs, Ps).
\end{lstlisting}
Conditions and faults are proved through the rules. Everything else, such as observations, thresholds, calculations and comparisons, is called directly and recorded as a leaf. \texttt{diagnose/2} collects every fault that can be proved, keeps the first proof of each, and applies the two closed-world conclusions: \texttt{normal} when no fault rule fires and every inlet is within the envelope, and \texttt{unexplained\_abnormality} when no fault rule fires but some inlet is not.

Negation as failure is used in two ways. It guards rules against known confounders. Rule r22, for example, concludes recirculation only if it cannot prove that the top inlet sensor is suspect. It also expresses the closed-world conclusions. Both uses stay within the conditions under which negation as failure is sound \parencite{clark1978}: the negated goal is ground when it is called.

Three groups of rules deserve comment:
\begin{itemize}
  \item \textbf{Backup rules} use independent evidence when the primary reading is missing. Rule r36 concludes a CRAH fan failure from row-wide overheating with unusually cold supply air, which does not need the fan power reading.
  \item \textbf{Sensor-fault rules} test a sensor against sensors that see the same air. Rule r43 goes further and accepts a supply-air sensor fault only if replacing its reading with the median middle inlet closes the coil energy balance (Equation~\eqref{eq:e41b}).
  \item \textbf{Physical-constraint rules} use the mixing model. Rule r46 flags a top inlet sensor that reads colder than the middle sensor, which Equation~\eqref{eq:mixing} rules out.
\end{itemize}

% ---------------------------------------------------------------------------
\section{Abductive Reasoning}
\label{sec:abductive}

Abduction implements Equations~\eqref{eq:abductive} and~\eqref{eq:parsimony} by parsimonious covering \parencite{peng1990}. The causal knowledge for each hypothesis says which symptoms it predicts, with what weight, and which of them it requires. Recirculation, for example, is described as:
\begin{lstlisting}
evidence(hot_air_recirculation(R), outlet_stratified(R),  1.0).
evidence(hot_air_recirculation(R), inlet_stratified(R),   0.8).
evidence(hot_air_recirculation(R), mid_inlet_deviates(R), 0.3).
requires(hot_air_recirculation(R), outlet_stratified(R)).
requires(hot_air_recirculation(R), inlet_stratified(R)).
masks(sensor_bias(t_inlet_top, R), inlet_stratified(R)).
\end{lstlisting}
\texttt{abduce/2} proceeds in four steps:
\begin{enumerate}
  \item It collects the clearly present symptoms that need explaining.
  \item It generates candidate hypotheses that explain at least one of them, and discards any that require a clearly absent symptom that no other candidate could mask.
  \item It searches subsets of the candidates in increasing size, up to four, and keeps every subset that covers all present symptoms and is consistent. A member may require a clearly absent symptom only if another member masks it. The first size that yields a solution gives the minimal explanations.
  \item It ranks the minimal explanations by summed confidence (Equation E56) and records a verdict for every candidate: part of an explanation, not needed, or rejected together with the absent symptom that ruled it out.
\end{enumerate}
Uncertain and unknown symptoms never need explaining, so a reading near the noise level cannot force an extra fault into the answer. Masking lets the module find a hidden fault: when the top inlet reads low over a recirculating rack, the explanation \{recirculation, top-sensor bias\} is consistent although recirculation alone is not.

% ---------------------------------------------------------------------------
\section{Consistency-Based Reasoning}
\label{sec:cbd}

Consistency-based diagnosis implements Equation~\eqref{eq:reiter}. Each analytical redundancy relation names the symptom that signals its violation and the components that must be healthy for it to hold. Each fault model names a relation that must be violated if a component is abnormal:
\begin{lstlisting}
arr(a_out_strat(R), outlet_stratified(R),
    [rack_seal(R), sensor(t_outlet_top, R), sensor(t_outlet, R)]).
arr(a_rack_power(R), rack_power_mismatch(R), [sensor(p_rack, R) | Phases]) :-
    component(R, rack),
    findall(sensor(i_branch, B), phase_of_rack(R, B), Phases).

fm(rack_seal(R), a_out_strat(R)).
fm(sensor(i_branch, B), a_rack_power(R)) :- phase_of_rack(R, B).
\end{lstlisting}
Every clearly violated relation yields a conflict set. \texttt{cbd\_diagnoses/2} then searches for minimal hitting sets \parencite{reiter1987}. It picks a conflict that the current candidate does not yet hit and branches on its members. When a candidate hits every conflict but breaks a fault model, because a component's predicted violation clearly did not happen, it branches on adding a component that could mask the effect. Candidates are capped at four components, and supersets of other diagnoses are removed.

Fault models are what let the module tell sensor faults from equipment faults \parencite{struss1989}. A branch overload produces the conflict \{branch, current sensor\}. A faulty current sensor would also break the rack power balance, so when that balance holds, the sensor is exonerated and only the overload remains. In the same way, a high top-inlet reading over a rack whose exhaust is not stratified exonerates the rack seal, so the diagnosis is a sensor fault.

% ---------------------------------------------------------------------------
\section{Explanation Facility}
\label{sec:explanation}

The explanation facility answers the question ``why was this diagnosis reached?''. \texttt{why/1} follows the workflow of Figure~\ref{fig:explanation}. For every reported fault that the rules prove, it reads the proof tree and reports the observed readings, the derived engineering quantities, the thresholds with their sources, the derived conditions, the rules fired with their provenance and the full reasoning trace. For a fault found only by abduction and consistency checking, it reports the graded supporting symptoms, each with its value, threshold, estimated noise, evidence strength and source. It then gives the maintenance plan and any plan interactions, the competing hypotheses with their confidence and the reason each was used, not needed or rejected, and the conflict sets from the consistency check.

\begin{figure}[htbp]
  \centering
  % Figure 6: explanation-generation workflow of why/1.
\begin{tikzpicture}[x=1cm, y=1cm]
  \node[out=5cm] (diag) at (4.5,0) {Final diagnosis and how it was decided};
  \node[decision=2.4cm] (dec) at (4.5,-2.0) {each label: proved by the rules?};
  \node[kb=4.6cm, align=left] (proof) at (0.9,-4.4)
    {\textbf{Deductive proof}\\
     \footnotesize observed readings\\
     \footnotesize derived engineering quantities\\
     \footnotesize thresholds with their sources\\
     \footnotesize derived conditions\\
     \footnotesize rules fired, with provenance\\
     \footnotesize full reasoning trace};
  \node[kb=4.8cm, align=left] (abd) at (8.2,-4.4)
    {\textbf{Abductive support}\\
     \footnotesize supporting symptoms, each with\\
     \footnotesize its value, threshold, noise,\\
     \footnotesize evidence strength and source};
  \node[out=5.8cm] (plan) at (4.5,-6.9) {Maintenance plan for each label\\ \footnotesize priority, ordered actions, rationale, source};
  \node[out=5.8cm] (inter) at (4.5,-8.4) {Plan interactions\\ \footnotesize which repair must come first};
  \node[reason=5.8cm] (comp) at (4.5,-9.9) {Competing hypotheses\\ \footnotesize confidence, verdict, reason for rejection};
  \node[reason=5.8cm] (cons) at (4.5,-11.4) {Consistency check\\ \footnotesize conflict sets, alternative minimal diagnoses};

  \draw[flow] (diag) -- (dec);
  \draw[flow] (dec.west) -| node[tag, pos=0.25, above] {yes} (proof.north);
  \draw[flow] (dec.east) -| node[tag, pos=0.25, above] {no} (abd.north);
  \draw[flow] (proof.south) |- (plan.west);
  \draw[flow] (abd.south) |- (plan.east);
  \draw[flow] (plan) -- (inter);
  \draw[flow] (inter) -- (comp);
  \draw[flow] (comp) -- (cons);
\end{tikzpicture}

  \caption{Explanation-generation workflow of the explanation facility.}
  \label{fig:explanation}
\end{figure}

Listing~\ref{lst:why} shows the explanation of a held-out test case in which recirculation at rack A04 is hidden by a top inlet sensor that reads low. The output is reproduced unedited. It shows each of the eight elements that the framework requires of an explanation: observed evidence, derived conditions, rules activated, intermediate reasoning, the diagnosis, competing hypotheses, the reasons for rejecting weaker ones and the recommended action.

\lstinputlisting[language={}, basicstyle=\ttfamily\scriptsize, caption={Unedited output of \texttt{why(h188)}: recirculation hidden by a low-reading sensor.}, label={lst:why}]{listings/why_h188.txt}

% ---------------------------------------------------------------------------
\section{Corrective-Maintenance Planner}
\label{sec:planner}

Every fault class has a plan in \texttt{kb/maintenance.pl}, with a priority, ordered actions, an engineering rationale and a source:
\begin{lstlisting}
maintenance(branch_overload, urgent,
    [ 'Identify the servers on the overloaded phase from the rack power records',
      'Rebalance server power supplies across the three phases',
      'Check for equipment added without a capacity review',
      'Confirm every phase is below 80% of its breaker rating' ],
    'A phase above its continuous rating can trip its breaker and drop every
     server on that circuit.',
    nec_continuous).
\end{lstlisting}
Each plan ends with a verification step that confirms the repair worked, such as re-measuring a temperature or a current. Figure~\ref{fig:maintenance} shows how plans are assembled. A diagnosed label is split into its fault class and entity, and the plan for the class is instantiated for the entity.

\begin{figure}[htbp]
  \centering
  % Figure 7: corrective-maintenance planning workflow.
\begin{tikzpicture}[x=1cm, y=1cm]
  \node[out=3.4cm] (lab) at (-0.5,0) {Diagnosed fault labels\\ \footnotesize e.g.\ a branch overload on phase A1 of rack A02};
  \node[reason=3.4cm] (map) at (4.4,0) {Split each label into fault class and entity};
  \node[kb=3.4cm] (kb) at (8.8,0) {Look up \texttt{maintenance/5}\\ \footnotesize priority, actions, rationale, source};
  \node[kb=3.4cm] (int) at (8.8,-2.2) {Check \texttt{\footnotesize plan\_interaction/3} for every pair of labels};
  \node[out=3.4cm] (plan) at (4.4,-2.2) {Ordered work plan\\ \footnotesize repairs that others depend on come first};
  \node[data=3.4cm] (eval) at (-0.5,-2.2) {PlanValidity\\ \footnotesize does each plan address a fault really present?};

  \draw[flow] (lab) -- (map);
  \draw[flow] (map) -- (kb);
  \draw[flow] (kb) -- (int);
  \draw[flow] (int) -- (plan);
  \draw[uses] (plan) -- node[tag, above] {evaluated} (eval);
\end{tikzpicture}

  \caption{Corrective-maintenance planning workflow.}
  \label{fig:maintenance}
\end{figure}

Table~\ref{tab:maintenance} lists the plan issued for every diagnosis. It is generated from the knowledge base, like the rule table.

% Generated by tests/export_latex.pl from the knowledge base. Do not edit by hand.
\begin{footnotesize}
\begin{longtable}{L{2.8cm}L{1.3cm}L{8.2cm}L{2.0cm}}
  \caption{Maintenance-action mapping, generated from \texttt{kb/maintenance.pl}: the corrective plan issued for each diagnosis, in order.}\label{tab:maintenance}\\
  \toprule
  Diagnosis & Priority & Ordered actions & Source \\
  \midrule
  \endfirsthead
  \toprule
  Diagnosis & Priority & Ordered actions & Source \\
  \midrule
  \endhead
  \texttt{crah\_\allowbreak{}fan\_\allowbreak{}failure} & urgent & 1.~Check the fan drive fault log and the motor current \newline 2.~Inspect the motor, belt or coupling, and bearings; repair or replace the failed part \newline 3.~Bring redundant cooling online or raise the output of neighbouring units during the repair \newline 4.~After repair, confirm fan power matches the affinity-law prediction and rack inlets fall below 27 C & \texttt{ashrae\_\allowbreak{}hb\_\allowbreak{}fans} \\
  \texttt{chw\_\allowbreak{}supply\_\allowbreak{}loss} & urgent & 1.~Confirm chilled-water plant status: pumps running and isolation valves open \newline 2.~Compare commanded and actual valve stroke; repair or replace a seized actuator \newline 3.~Vent trapped air from the coil \newline 4.~Confirm water flow matches the valve position and supply air is back at setpoint & \texttt{a4\_\allowbreak{}chw\_\allowbreak{}flow} \\
  \texttt{hot\_\allowbreak{}air\_\allowbreak{}recirculation} & high & 1.~Inspect the rack for missing blanking panels and unsealed cable openings \newline 2.~Check containment doors and roof panels for gaps \newline 3.~Install blanking panels and seal the openings \newline 4.~Re-measure top and middle inlet temperatures to confirm the difference is below 3 K & \texttt{eu\_\allowbreak{}coc\_\allowbreak{}dc} \\
  \texttt{normal} & routine & 1.~No corrective action. Continue routine monitoring. & \texttt{ashrae\_\allowbreak{}tc99\_\allowbreak{}2021} \\
  \texttt{unexplained\_\allowbreak{}abnormality} & high & 1.~Verify the affected sensors before acting \newline 2.~Escalate to a technician for manual investigation & \texttt{closed\_\allowbreak{}world} \\
  \texttt{filter\_\allowbreak{}clogging} & high & 1.~Replace the CRAH air filters \newline 2.~Check the filter differential pressure gauge against a reference \newline 3.~Review the filter replacement interval and recent dust sources such as building work \newline 4.~After replacement, confirm filter pressure drop and fan speed are back to normal & \texttt{a8\_\allowbreak{}filter} \\
  \texttt{server\_\allowbreak{}fan\_\allowbreak{}failure} & high & 1.~Check the server management controller event log for the failed fan module \newline 2.~Migrate or reduce the workload on the server until it is repaired \newline 3.~Hot-swap the failed fan module \newline 4.~Confirm CPU temperature and clock frequency return to normal & \texttt{a9\_\allowbreak{}server\_\allowbreak{}fan} \\
  \texttt{branch\_\allowbreak{}overload} & urgent & 1.~Identify the servers on the overloaded phase from the rack power records \newline 2.~Rebalance server power supplies across the three phases \newline 3.~Check for equipment added without a capacity review \newline 4.~Confirm every phase is below 80\% of its breaker rating & \texttt{nec\_\allowbreak{}continuous} \\
  \texttt{battery\_\allowbreak{}degradation} & high & 1.~Confirm with a capacity discharge test as IEEE Std 1188 recommends \newline 2.~Check battery room temperature and its cooling \newline 3.~Replace weak blocks or the whole string, depending on the test result \newline 4.~Record a new resistance baseline after replacement & \texttt{a11\_\allowbreak{}battery} \\
  \texttt{ups\_\allowbreak{}on\_\allowbreak{}bypass} & urgent & 1.~Read the UPS alarm log to find why it transferred to bypass \newline 2.~Call the UPS service engineer to repair the inverter or rectifier fault \newline 3.~Postpone any upstream electrical work while the load has no battery protection \newline 4.~Return the UPS to double conversion after the repair and confirm normal efficiency & \texttt{a12\_\allowbreak{}ups} \\
  \texttt{sensor\_\allowbreak{}bias} & high & 1.~Check the sensor against a calibrated reference instrument \newline 2.~Check the sensor mounting and position \newline 3.~Recalibrate or replace the sensor \newline 4.~Do not act on readings from this sensor until it is verified & \texttt{a13\_\allowbreak{}neighbour} \\
  \texttt{telemetry\_\allowbreak{}dropout} & high & 1.~Check the monitoring agent, controller and network link for this telemetry group \newline 2.~Restore the data feed \newline 3.~Treat the equipment status as unknown until fresh data arrives \newline 4.~Inspect critical equipment on site if the feed cannot be restored quickly & \texttt{a14\_\allowbreak{}telemetry} \\
  \bottomrule
\end{longtable}
\end{footnotesize}


\paragraph{Plan interactions.} When two diagnoses interact, their plans must be ordered. The verification step of the recirculation plan re-measures the rack's top inlet temperature. If the same rack's top inlet sensor is also diagnosed as biased, that step would rely on a faulty sensor. \texttt{plan\_interaction/3} records such dependencies, and the planner then says to repair the sensor first, as Listing~\ref{lst:why} shows. Similar interactions cover the supply-air sensor, whose reading the cooling-unit plans rely on, and stale telemetry, which leaves any plan for the affected equipment unverifiable.

\paragraph{Plan validity.} Whether each issued plan addresses a fault that is really present is measured as PlanValidity in Chapter~\ref{ch:dataset_and_evaluation}. A plan for a fault that is not present wastes maintenance effort. A missing plan for a fault that is present leaves it unrepaired.

\section*{Generated tables}
\addcontentsline{toc}{section}{Generated tables}

Tables~\ref{tab:thresholds} and~\ref{tab:provenance} are generated from the knowledge base by \texttt{tests/export\_latex.pl}, like Table~\ref{tab:rules}.

% Generated by tests/export_latex.pl from the knowledge base. Do not edit by hand.
\begin{small}
\begin{longtable}{L{4.0cm}L{1.4cm}L{1.5cm}L{3.7cm}L{2.4cm}}
  \caption{Thresholds, generated from \texttt{kb/thresholds.pl}.}\label{tab:thresholds}\\
  \toprule
  Threshold & Value & Unit & Source & Source type \\
  \midrule
  \endfirsthead
  \toprule
  Threshold & Value & Unit & Source & Source type \\
  \midrule
  \endhead
  \texttt{t\_\allowbreak{}inlet\_\allowbreak{}rec\_\allowbreak{}max} & 27.0 & degC & \texttt{ashrae\_\allowbreak{}tc99\_\allowbreak{}2021} & standard \\
  \texttt{t\_\allowbreak{}inlet\_\allowbreak{}rec\_\allowbreak{}min} & 18.0 & degC & \texttt{ashrae\_\allowbreak{}tc99\_\allowbreak{}2021} & standard \\
  \texttt{t\_\allowbreak{}supply\_\allowbreak{}setpoint} & 20.0 & degC & \texttt{a1\_\allowbreak{}supply\_\allowbreak{}setpoint} & assumption \\
  \texttt{t\_\allowbreak{}supply\_\allowbreak{}band} & 2.0 & kelvin & \texttt{a1\_\allowbreak{}supply\_\allowbreak{}setpoint} & assumption \\
  \texttt{row\_\allowbreak{}hot\_\allowbreak{}fraction} & 0.5 & ratio & \texttt{a2\_\allowbreak{}row\_\allowbreak{}fraction} & assumption \\
  \texttt{valve\_\allowbreak{}open\_\allowbreak{}min} & 50.0 & percent & \texttt{a3\_\allowbreak{}valve\_\allowbreak{}open} & assumption \\
  \texttt{chw\_\allowbreak{}flow\_\allowbreak{}ratio\_\allowbreak{}min} & 0.3 & ratio & \texttt{a4\_\allowbreak{}chw\_\allowbreak{}flow} & assumption \\
  \texttt{strat\_\allowbreak{}max} & 3.0 & kelvin & \texttt{a5\_\allowbreak{}stratification} & assumption \\
  \texttt{fan\_\allowbreak{}ratio\_\allowbreak{}min} & 0.5 & ratio & \texttt{a6\_\allowbreak{}fan\_\allowbreak{}ratio} & assumption \\
  \texttt{strat\_\allowbreak{}out\_\allowbreak{}max} & 3.5 & kelvin & \texttt{a7\_\allowbreak{}outlet\_\allowbreak{}strat} & assumption \\
  \texttt{filter\_\allowbreak{}loading\_\allowbreak{}max} & 1.7 & ratio & \texttt{a8\_\allowbreak{}filter} & assumption \\
  \texttt{server\_\allowbreak{}fan\_\allowbreak{}min} & 1000.0 & rpm & \texttt{a9\_\allowbreak{}server\_\allowbreak{}fan} & assumption \\
  \texttt{cpu\_\allowbreak{}hot\_\allowbreak{}min} & 85.0 & degC & \texttt{a10\_\allowbreak{}cpu} & assumption \\
  \texttt{branch\_\allowbreak{}loading\_\allowbreak{}max} & 0.8 & ratio & \texttt{nec\_\allowbreak{}continuous} & standard \\
  \texttt{cell\_\allowbreak{}r\_\allowbreak{}ratio\_\allowbreak{}max} & 1.3 & ratio & \texttt{a11\_\allowbreak{}battery} & assumption \\
  \texttt{ups\_\allowbreak{}eff\_\allowbreak{}bypass\_\allowbreak{}min} & 0.98 & ratio & \texttt{a12\_\allowbreak{}ups} & assumption \\
  \texttt{neighbour\_\allowbreak{}dev\_\allowbreak{}max} & 2.0 & kelvin & \texttt{a13\_\allowbreak{}neighbour} & assumption \\
  \texttt{balance\_\allowbreak{}resid\_\allowbreak{}max} & 0.2 & ratio & \texttt{eq40\_\allowbreak{}uncertainty} & method \\
  \texttt{sample\_\allowbreak{}age\_\allowbreak{}max} & 60.0 & s & \texttt{a14\_\allowbreak{}telemetry} & assumption \\
  \texttt{t\_\allowbreak{}supply\_\allowbreak{}hot\_\allowbreak{}limit} & 22.0 & degC & \texttt{a1\_\allowbreak{}supply\_\allowbreak{}setpoint} & assumption \\
  \texttt{fan\_\allowbreak{}ratio\_\allowbreak{}reduced} & 0.85 & ratio & \texttt{a17\_\allowbreak{}fan\_\allowbreak{}reduced} & assumption \\
  \texttt{t\_\allowbreak{}supply\_\allowbreak{}cold\_\allowbreak{}limit} & 18.0 & degC & \texttt{a1\_\allowbreak{}supply\_\allowbreak{}setpoint} & assumption \\
  \texttt{power\_\allowbreak{}resid\_\allowbreak{}max} & 0.1 & ratio & \texttt{a15\_\allowbreak{}power\_\allowbreak{}balance} & assumption \\
  \texttt{sensor\_\allowbreak{}noise\_\allowbreak{}floor} & 0.15 & kelvin & \texttt{a16\_\allowbreak{}noise} & assumption \\
  \bottomrule
\end{longtable}
\end{small}

% Generated by tests/export_latex.pl from the knowledge base. Do not edit by hand.
\begin{small}
\begin{longtable}{L{3.4cm}L{2.4cm}L{8.2cm}}
  \caption{Provenance of rules and thresholds, generated from \texttt{kb/sources.pl}.}\label{tab:provenance}\\
  \toprule
  Source & Type & Citation or justification \\
  \midrule
  \endfirsthead
  \toprule
  Source & Type & Citation or justification \\
  \midrule
  \endhead
  \texttt{ashrae\_\allowbreak{}tc99\_\allowbreak{}2021} & standard & ASHRAE TC 9.9 (2021). Thermal guidelines for data processing environments (5th ed.). ASHRAE. \\
  \texttt{ashrae\_\allowbreak{}hb\_\allowbreak{}fans} & literature & ASHRAE (2020). 2020 ASHRAE Handbook - HVAC Systems and Equipment, fans chapter. Fan affinity laws. \\
  \texttt{eu\_\allowbreak{}coc\_\allowbreak{}dc} & standard & Acton et al. (2018). 2018 Best Practice Guidelines for the EU Code of Conduct on Data Centre Energy Efficiency, v9.1.0 (European Commission JRC), practice 5.1.4. \\
  \texttt{a1\_\allowbreak{}supply\_\allowbreak{}setpoint} & assumption & Design assumption: CRAH supply-air setpoint 20 C with a 2 K tolerance band. \\
  \texttt{a2\_\allowbreak{}row\_\allowbreak{}fraction} & assumption & Engineering assumption: overheating at half or more of the racks indicates a row-level cause. \\
  \texttt{a3\_\allowbreak{}valve\_\allowbreak{}open} & assumption & Engineering assumption: a valve open 50\% or more should produce a measurable water temperature rise. \\
  \texttt{a4\_\allowbreak{}chw\_\allowbreak{}flow} & assumption & Engineering assumption: measured water flow below 30\% of the valve-characteristic prediction means chilled-water supply is lost. \\
  \texttt{a5\_\allowbreak{}stratification} & assumption & Engineering assumption: a top-to-middle inlet difference above 3 K indicates hot-air recirculation. \\
  \texttt{a6\_\allowbreak{}fan\_\allowbreak{}ratio} & assumption & Engineering assumption: fan power below 50\% of the affinity-law prediction means the fan is not moving air. \\
  \texttt{closed\_\allowbreak{}world} & method & Closed-world assumption: normal operation is concluded when no fault rule fires and all inlets are within the envelope. \\
  \texttt{nec\_\allowbreak{}continuous} & standard & NFPA (2023). NFPA 70 National Electrical Code, continuous-load rule: 80\% of breaker rating. Local codes may differ. \\
  \texttt{a7\_\allowbreak{}outlet\_\allowbreak{}strat} & assumption & Engineering assumption: a top-to-bottom outlet difference above 3.5 K confirms recirculation. Normal racks stay below 3 K. \\
  \texttt{a8\_\allowbreak{}filter} & assumption & Engineering assumption: filter resistance 1.7 times its clean value calls for replacement (check against manufacturer final-resistance data). \\
  \texttt{a9\_\allowbreak{}server\_\allowbreak{}fan} & assumption & Engineering assumption: a server fan below 1000 rpm has failed. The controller never runs it below about 5400 rpm. \\
  \texttt{a10\_\allowbreak{}cpu} & assumption & Engineering assumption: CPU temperature of 85 C or more, 10 K below the 95 C throttling limit, confirms lost server cooling. \\
  \texttt{a11\_\allowbreak{}battery} & assumption & Engineering assumption: cell internal resistance 30\% above its baseline indicates degradation (check against IEEE Std 1188 and manufacturer guidance). \\
  \texttt{a12\_\allowbreak{}ups} & assumption & Engineering assumption: UPS efficiency of 0.98 or more is only reached on static bypass. Double conversion runs at about 0.95. \\
  \texttt{a13\_\allowbreak{}neighbour} & assumption & Engineering assumption: a temperature sensor more than 2 K from the median of sensors that see the same air is suspect (Stage 4 Eq. 41). \\
  \texttt{eq40\_\allowbreak{}uncertainty} & method & Residual threshold from uncertainty propagation, Stage 4 Eqs. 39 and 40, with a coverage factor of about 3. \\
  \texttt{a14\_\allowbreak{}telemetry} & assumption & Engineering assumption: a sample older than 60 s is stale, given a polling interval of about 30 s. \\
  \texttt{eq18\_\allowbreak{}mixing} & method & Physical constraint from the mixing model, Stage 4 Eq. 18: the top inlet cannot be colder than the middle inlet. \\
  \texttt{a15\_\allowbreak{}power\_\allowbreak{}balance} & assumption & Engineering assumption: metered rack power and the sum of its six phase currents agree within 10\% when both meters are healthy (Stage 4 Eq. 42). \\
  \texttt{a16\_\allowbreak{}noise} & assumption & Engineering assumption: temperature-sensor noise is at least 0.15 K (typical datasheet accuracy). Higher noise is estimated from the spread of middle inlets, which all see the same supply air. \\
  \texttt{a17\_\allowbreak{}fan\_\allowbreak{}reduced} & assumption & Engineering assumption: fan power below 85\% of the affinity-law prediction shows a partial failure such as belt slip. Healthy fans stay within 5\%. \\
  \bottomrule
\end{longtable}
\end{small}

